\begin{proposition}[7.9.8]
\label{IV.2.s7.9.8-fr}
Sean $A$ un anillo local noetheriano, $k$ su cuerpo residual, $B$ un anillo
local noetheriano tal que $\operatorname{Spec}(B)\in\boldsymbol{C}$,
$\varphi:A\to B$ un homomorfismo local que haga de $B$ un $A$-módulo plano;
pongamos $Y=\operatorname{Spec}(A)$, $X=\operatorname{Spec}(B)$ y sea
$f:X\to Y$ el morfismo correspondiente a $\varphi$. Se supone que:

$1^\circ$ la propiedad $\boldsymbol P(B\otimes_A k,k)$ sea verdadera;

$2^\circ$ para todo morfismo finito $Y_1\to Y$, se pueda resolver
$(Y_1)_{\mathrm{red}}$.

Entonces, para todo $y\in Y$, la propiedad
$\boldsymbol P(f^{-1}(y),\boldsymbol{k}(y))$ es verdadera.
\end{proposition}

\begin{proof}
Observemos que la condición $2^\circ$ de (7.9.8) sigue satisfaciéndose si
se sustituye $Y$ por el espectro del anillo local en un ideal maximal de una
$A$-álgebra finita $A'$ (7.9.2); por otra parte, en virtud de la condición
$1^\circ$ de (7.9.7), el espectro de una localización de
$B\otimes_AA'$ pertenece también a $\boldsymbol{C}$. El lema (7.3.16.2) muestra
entonces (como en la parte I) del razonamiento de (7.5.2)) que basta
demostrar que, si $Y$ es íntegro y $y$ es el punto genérico de $Y$, los
anillos locales de la fibra $f^{-1}(y)$ satisfacen $\boldsymbol R$.

Así pues, por hipótesis existe un preesquema \emph{regular} $Y'$ y un morfismo
propio y birracional $g:Y'\to Y$; como $Y'$ es noetheriano y localmente
íntegro, la hipótesis de que $g$ sea birracional implica que $Y'$ es íntegro;
sea $X'=X\times_YY'$, de modo que se tiene un diagrama conmutativo

\begin{env}[7.9.8.1]
\label{IV.2.s7.9.8.1-fr}
\[
\xymatrix{
X\ar[d]_f & X'\ar[l]_{g'}\ar[d]^{f'}\\
Y & Y'\ar[l]^g
}
\]
\end{env}

en el que $f'$ y $g'$ son las proyecciones canónicas. Teniendo en cuenta la condición
$2^\circ$ de (7.9.7), el mismo razonamiento que al comienzo de (7.9.3) muestra
que basta demostrar

\oldpage[IV-2]{222}

que se tiene $U_{\mathbf{R}}(X')=X'$; a su vez, denotando por $a$ el punto
cerrado de $X$, basta demostrar que $g'^{-1}(a)\subset U_{\mathbf{R}}(X')$. Ahora bien, sea
$x'\in g'^{-1}(a)$ y pongamos $z'=f'(x')$, de modo que $b=g(z')$ es el
punto cerrado de $Y$. Se tiene, por tanto, $f'^{-1}(z')=f^{-1}(b)\otimes_k \boldsymbol{k}(z')$,
y como $\boldsymbol{k}(z')$ es una extensión \emph{de tipo finito} de $k$ (pues $g$ es
propio y, por tanto, de tipo finito), la hipótesis de que $\boldsymbol P(f^{-1}(b),k)$
sea verdadera implica, en virtud de la condición $4^\circ$ de (7.9.7), que también
lo es $\boldsymbol P(f'^{-1}(z'),\boldsymbol{k}(z'))$. Pero como
$\mathcal O_{z'}$ es un anillo regular y
$\operatorname{Spec}(\mathcal O_{x'})\in\boldsymbol{C}$ en virtud de la
condición $1^\circ$ de (7.9.7), el lema (7.5.1.1) prueba que
$\boldsymbol R(\mathcal O_{X',x'})$ es verdadera, con lo que termina la demostración.
\end{proof}
